\subsection*{מטרות השיעור}

\begin{enumerate}
    \item פתרון המשוואה הרדיאלית לחלקיק בפוטנציאל קולוני $V(r) = -\frac{Ze^2}{4\pi\epsilon_0 r}$
    \item הפרדה מלאה למשתנים: חלק רדיאלי $R(r)$ וחלק זוויתי $Y_{lm}(\theta,\phi)$
    \item תנאי שפה על הפונקציה הרדיאלית והרמיטיות אופרטור התנע $\hat{p}_r$
    \item קוונטיזציה: מציאת ספקטרום האנרגיה של אטום המימן
    \item פולינומי לאגר והפונקציות העצמיות המלאות
    \item הבנת הניוון $n^2$ וקשר לסימטריה נסתרת (וקטור \LR{Runge-Lenz})
    \item סקלות אופייניות: רדיוס בוהר, קבוע המבנה הדק, אנרגיית ריידברג
    \item משמעות פיזיקלית: קווים ספקטרליים, מעברים מותרים ואסורים
\end{enumerate}

\subsection*{תזכורת — הבעיה הדו-גופית והמילטוניאן}

\begin{summarybox}{מה הגענו עד כה}
\textbf{בשיעור הקודם:}
\begin{itemize}
    \item הפרדנו את הבעיה הדו-גופית למרכז מסה + תנועה יחסית
    \item המסה המצומצמת: $\mu = \frac{M_1 M_2}{M_1 + M_2} \approx m_e$ (עבור אטום מימן)
    \item המילטוניאן: $\hat{H} = -\frac{\hbar^2}{2\mu}\nabla^2 + V(r)$
    \item בפוטנציאל מרכזי: $\hat{H} = -\frac{\hbar^2}{2\mu r^2}\frac{d}{dr}\left(r^2\frac{d}{dr}\right) + \frac{\hat{L}^2}{2\mu r^2} + V(r)$
\end{itemize}

\textbf{היום:} נפתור את המשוואה הרדיאלית במפורש עבור פוטנציאל קולוני!
\end{summarybox}

\subsection{המשוואה הרדיאלית והפרדת משתנים}

\subsubsection{ההמילטוניאן בקואורדינטות כדוריות}

עבור פוטנציאל מרכזי $V(r)$, ההמילטוניאן:
\[
\hat{H} = -\frac{\hbar^2}{2\mu}\left[\frac{\partial^2}{\partial r^2} + \frac{2}{r}\frac{\partial}{\partial r}\right] + \frac{\hat{L}^2}{2\mu r^2} + V(r)
\]

\begin{importantbox}{הפרדת משתנים}
מחפשים פתרון מהצורה:
\[
\psi(r,\theta,\phi) = R(r) \cdot Y_{lm}(\theta,\phi)
\]

כאשר $Y_{lm}$ הן הרמוניות ספריות המקיימות:
\begin{align*}
\hat{L}^2 Y_{lm} &= \hbar^2 l(l+1) Y_{lm} \\
\hat{L}_z Y_{lm} &= \hbar m Y_{lm}
\end{align*}

הצבה ב-$\hat{H}\psi = E\psi$ נותנת משוואה רדיאלית עבור $R(r)$!
\end{importantbox}

\subsection{תנאי שפה והרמיטיות אופרטור התנע הרדיאלי}

\subsubsection{הגדרת אופרטור התנע הרדיאלי}

\begin{definitionbox}{אופרטור התנע הרדיאלי $\hat{p}_r$}
האופרטור הנאיבי $\hat{p}_r = -i\hbar \frac{\partial}{\partial r}$ **אינו הרמיטי**!

האופרטור ההרמיטי הנכון:
\[
\hat{p}_r = -i\hbar\left(\frac{\partial}{\partial r} + \frac{1}{r}\right)
\]

\textbf{למה?} כדי לשמור על הרמיטיות בקואורדינטות כדוריות עם יעקוביאן $r^2$.
\end{definitionbox}

\textbf{דרישת הרמיטיות ותנאי שפה}
אופרטור $\hat{A}$ הרמיטי אם: $\braket{\psi|\hat{A}|\phi} = \braket{\hat{A}\psi|\phi}^*$

עבור $\hat{p}_r$, דרישת הרמיטיות:
\[
0 = \int_0^\infty dr \, r^2 \left[\psi^*\hat{p}_r\phi - (\hat{p}_r\psi)^*\phi\right]
\]

\textbf{שלב 1:} אינטגרציה בחלקים:
\[
= -i\hbar \int_0^\infty dr \left[\frac{d}{dr}(r^2\psi^*\psi)\right] = -i\hbar \left[r^2\psi^*\psi\right]_0^\infty
\]

\textbf{שלב 2:} תנאי שפה:
\begin{itemize}
    \item ב-$r \to \infty$: $\psi \to 0$ (מצב קשור)
    \item ב-$r \to 0$: $r\psi \to 0$
\end{itemize}

\textbf{מסקנה:} $R(r)$ חייבת לקיים $\lim_{r \to 0} r \cdot R(r) = 0$

\subsubsection{$\hat{p}_r^2$ והקשר ללפלסיאן}

\begin{theorembox}{קשר בין $\hat{p}_r^2$ ללפלסיאן}
\[
\hat{p}_r^2 = -\hbar^2\left(\frac{\partial^2}{\partial r^2} + \frac{2}{r}\frac{\partial}{\partial r}\right)
\]

לכן:
\[
\hat{H} = \frac{\hat{p}_r^2}{2\mu} + \frac{\hat{L}^2}{2\mu r^2} + V(r)
\]

זהו בדיוק המילטוניאן של בעיה חד-ממדית רדיאלית עם "פוטנציאל אפקטיבי"!
\end{theorembox}

\begin{notebox}
\textbf{הערה חשובה:} למרות ש-$\hat{p}_r$ הרמיטי, **אין לו ערכי תצפית** (ספקטרום)!

הסיבה: אף פונקציה לא מקיימת בו-זמנית $r\psi(0) = 0$ ו-$\psi(\infty) = 0$ ותהיה פונקציה עצמית של $\hat{p}_r$.

זו "חיה מוזרה" — אופרטור הרמיטי שאין לו ספקטרום פיזיקלי!
\end{notebox}

\textbf{ההחלפה $u(r) = r \cdot R(r)$}

\begin{importantbox}{טריק מתמטי חשוב}
נגדיר: $u(r) = r \cdot R(r)$

\textbf{יתרונות:}
\begin{enumerate}
    \item תנאי השפה: $u(0) = 0$ ו-$u(\infty) = 0$
    \item המשוואה מתפשטת: $\frac{d^2u}{dr^2} + \frac{2}{r}\frac{du}{dr} = \frac{1}{r}\frac{d^2(ru)}{dr^2}$
\end{enumerate}

\textbf{המשוואה הרדיאלית עבור $u(r)$:}
\[
-\frac{\hbar^2}{2\mu}\frac{d^2u}{dr^2} + \left[\frac{\hbar^2 l(l+1)}{2\mu r^2} + V(r)\right]u = E u
\]

זו משוואת שרדינגר חד-ממדית עם פוטנציאל אפקטיבי:
\[
V_{\text{eff}}(r) = V(r) + \frac{\hbar^2 l(l+1)}{2\mu r^2}
\]
\end{importantbox}

\subsection{פוטנציאל קולוני ואטום המימן}

\subsubsection{הפוטנציאל ויחידות}

\begin{definitionbox}{פוטנציאל קולוני}
\[
V(r) = -\frac{Ze^2}{4\pi\epsilon_0 r}
\]

כאשר:
\begin{itemize}
    \item $Z$ — מספר פרוטונים בגרעין ($Z=1$ למימן, $Z=2$ להליום, וכו')
    \item $e$ — מטען האלקטרון
    \item עבור מימן: $\mu \approx m_e$ (מסת האלקטרון)
\end{itemize}
\end{definitionbox}

\subsubsection{סקלות אופייניות}

\begin{theorembox}{סקלות האופייניות לאטום מימן}

\textbf{1. רדיוס בוהר $a_0$:}
\[
a_0 = \frac{4\pi\epsilon_0\hbar^2}{\mu e^2} \approx 0.529 \times 10^{-10} \text{ m} = 0.529 \, \text{\AA}
\]

זהו סדר הגודל של אטום המימן!

\textbf{2. קבוע המבנה הדק $\alpha$:}
\[
\alpha = \frac{e^2}{4\pi\epsilon_0\hbar c} \approx \frac{1}{137.036}
\]

פרמטר חסר ממדים שקובע את עוצמת האינטראקציה האלקטרומגנטית.

\textbf{3. אנרגיית ריידברג $R_\infty$:}
\[
R_\infty = \frac{\alpha^2 m_e c^2}{2} = \frac{e^2}{8\pi\epsilon_0 a_0} \approx 13.6 \, \text{eV}
\]

אנרגיית הקישור של אלקטרון ברמת היסוד של מימן!
\end{theorembox}

\begin{examplebox}{השוואת סקלות}
\textbf{גודל האטום \LR{vs.} גודל הגרעין:}
\begin{itemize}
    \item רדיוס בוהר: $a_0 \sim 10^{-10}$ m
    \item רדיוס הגרעין: $r_{\text{nucleus}} \sim 10^{-15}$ m
    \item יחס: $\frac{a_0}{r_{\text{nucleus}}} \sim 10^5$
\end{itemize}

\textbf{מסקנה:} הגרעין נראה כמעט כנקודתי מנקודת מבטו של האלקטרון!

\textbf{מסה מצומצמת:}
\[
\mu = \frac{m_e M}{m_e + M} \approx m_e\left(1 - \frac{m_e}{M}\right) \approx m_e \cdot (1 - 5\times10^{-4})
\]

התיקון קטן מאוד אבל ניתן למדידה בספקטרוסקופיה דיוק גבוה!
\end{examplebox}

\subsection{עבור למשתנים חסרי ממד}

\textbf{נרמול המשוואה}
\textbf{משתנים חדשים:}
\begin{align*}
x &= \frac{r}{a_0} \quad \text{(אורך חסר ממד)} \\
\varepsilon &= \frac{E}{\alpha^2 m_e c^2} \quad \text{(אנרגיה חסרת ממד)}
\end{align*}

\textbf{המשוואה הרדיאלית הופכת ל:}
\[
-\frac{1}{2}\frac{d^2u}{dx^2} + \left[\frac{l(l+1)}{2x^2} - \frac{Z}{x}\right]u = \varepsilon u
\]

נקייה ופשוטה — אין פרמטרים עם ממדים!

\subsection{ניתוח אסימפטוטי והצבה}

\subsubsection{גבול $x \to 0$}

\textbf{התנהגות ליד הראשית}
כאשר $x \to 0$, האיבר הדומיננטי:
\[
-\frac{1}{2}\frac{d^2u}{dx^2} + \frac{l(l+1)}{2x^2}u \approx 0
\]

\textbf{פתרון:} $u(x) \sim x^\lambda$

הצבה: $\lambda(\lambda-1) = l(l+1)$

\textbf{פתרונות:}
\begin{itemize}
    \item $\lambda = l+1$ ✓ (רגולרי)
    \item $\lambda = -l$ ✗ (סינגולרי ב-$x=0$)
\end{itemize}

\textbf{מסקנה:} $u(x) \sim x^{l+1}$ כאשר $x \to 0$


\subsubsection{x שואף לאינסוף}

\textbf{התנהגות באינסוף}
עבור מצבים קשורים ($\varepsilon < 0$), נגדיר:
\[
\kappa = \sqrt{-2\varepsilon} > 0
\]

באינסוף:
\[
-\frac{1}{2}\frac{d^2u}{dx^2} \approx \varepsilon u = -\frac{\kappa^2}{2}u
\]

\textbf{פתרון:} $u(x) \sim e^{-\kappa x}$ (דעיכה אקספוננציאלית)


\subsubsection{ההצבה הכללית}

\begin{importantbox}{אנזץ (\LR{Ansatz}) לפתרון}
\[
u(x) = x^{l+1} e^{-\kappa x} P(x)
\]

כאשר $P(x)$ פולינום שיש לקבוע.

\textbf{הרעיון:}
\begin{itemize}
    \item $x^{l+1}$ מטפל בהתנהגות ליד $x=0$
    \item $e^{-\kappa x}$ מטפל בהתנהגות ב-$x \to \infty$
    \item $P(x)$ מטפל ב"אמצע"
\end{itemize}
\end{importantbox}

\subsection{מעבר למשתנה ומשוואת לאגר}

\textbf{משתנה חדש $\rho$}
נגדיר:
\begin{align*}
n &= \frac{Z}{\kappa} = \frac{Z}{\sqrt{-2\varepsilon}} \\
\rho &= \frac{2Z}{n}x = 2\kappa x
\end{align*}

\textbf{הצבה:} $u(\rho) = \rho^{l+1} e^{-\rho/2} F(\rho)$

המשוואה ל-$F(\rho)$:
\[
\rho F'' + (2l+2-\rho)F' + (n-l-1)F = 0
\]

זוהי **משוואת לאגר המוכללת**!


\subsection{פולינומי לאגר וקוונטיזציה}

\subsubsection{פיתוח בטור}

\begin{definitionbox}{פיתוח טורים}
\[
F(\rho) = \sum_{q=0}^\infty a_q \rho^q
\]

הצבה במשוואה נותנת יחס רקורסיה:
\[
a_{q+1} = \frac{q+l+1-n}{(q+1)(2l+2+q)} a_q
\]
\end{definitionbox}

\begin{importantbox}{תנאי הקוונטיזציה}
\textbf{בעיה:} כאשר $q \to \infty$, $a_{q+1}/a_q \to 1/q$

זה אומר $F(\rho) \sim e^\rho$ — מתפוצץ באינסוף! \xmark

\textbf{פתרון:} הטור חייב להסתיים!

\textbf{דרישה:} קיים $q_{\max}$ כך ש:
\[
n = q_{\max} + l + 1
\]

\textbf{מכאן:}
\begin{itemize}
    \item $n \in \mathbb{N}$, $n \geq 1$ — **המספר הקוונטי הראשי**
    \item $l = 0, 1, 2, \ldots, n-1$ — **תנע זוויתי מקסימלי**
    \item $m = -l, \ldots, +l$
\end{itemize}
\end{importantbox}

\subsection{ספקטרום האנרגיה של אטום המימן}

\begin{theorembox}{נוסחת בוהר לאנרגיה}
מהקוונטיזציה $n = Z/\kappa$:
\[
\kappa = \frac{Z}{n} \quad \Rightarrow \quad \varepsilon = -\frac{\kappa^2}{2} = -\frac{Z^2}{2n^2}
\]

\textbf{האנרגיה עם יחידות:}
\[
E_n = -\frac{Z^2}{n^2} \cdot \frac{\alpha^2 m_e c^2}{2} = -\frac{Z^2}{n^2} R_\infty
\]

עבור מימן ($Z=1$):
\[
\boxed{E_n = -\frac{13.6 \, \text{eV}}{n^2}, \quad n = 1, 2, 3, \ldots}
\]

זוהי **נוסחת בוהר** — התוצאה ההיסטורית המפורסמת!
\end{theorembox}

\begin{examplebox}{רמות האנרגיה הראשונות}
\begin{align*}
E_1 &= -13.6 \, \text{eV} \quad \text{(רמת יסוד)} \\
E_2 &= -3.4 \, \text{eV} \\
E_3 &= -1.51 \, \text{eV} \\
E_\infty &= 0 \quad \text{(סף יינון)}
\end{align*}

ההפרש $E_\infty - E_1 = 13.6$ eV הוא אנרגיית היינון של אטום מימן.
\end{examplebox}

\subsection{הפונקציות העצמיות והניוון}

\subsubsection{הפונקציה הרדיאלית}

\begin{definitionbox}{פולינומי לאגר המוכללים}
\[
R_{nl}(r) = \sqrt{\left(\frac{2Z}{na_0}\right)^3 \frac{(n-l-1)!}{2n[(n+l)!]}} \cdot e^{-Zr/na_0} \left(\frac{2Zr}{na_0}\right)^l L_{n-l-1}^{2l+1}\left(\frac{2Zr}{na_0}\right)
\]

כאשר $L_q^{2l+1}(\rho)$ הם פולינומי לאגר המוכללים.
\end{definitionbox}

\subsubsection{הפונקציה המלאה}

\begin{importantbox}{פונקציות עצמיות של אטום המימן}
\[
\psi_{nlm}(r,\theta,\phi) = R_{nl}(r) \cdot Y_{lm}(\theta,\phi)
\]

\textbf{מספרים קוונטיים:}
\begin{itemize}
    \item $n = 1, 2, 3, \ldots$ — מספר קוונטי ראשי (קובע אנרגיה)
    \item $l = 0, 1, \ldots, n-1$ — תנע זוויתי אורביטלי
    \item $m = -l, -l+1, \ldots, l$ — הקרנה על ציר $z$
\end{itemize}
\end{importantbox}

\subsection{הניוון והסימטריה הנסתרת}

\subsubsection{ספירת מצבים}

\begin{examplebox}{ניוון רמות האנרגיה}

\textbf{רמת יסוד ($n=1$):}
\begin{itemize}
    \item $l=0$, $m=0$ $\Rightarrow$ 1 מצב
\end{itemize}

\textbf{רמה ראשונה מעוררת ($n=2$):}
\begin{itemize}
    \item $l=0$, $m=0$ $\Rightarrow$ 1 מצב
    \item $l=1$, $m=-1,0,+1$ $\Rightarrow$ 3 מצבים
    \item סה"כ: 4 מצבים
\end{itemize}

\textbf{רמה שנייה מעוררת ($n=3$):}
\begin{itemize}
    \item $l=0$: 1 מצב
    \item $l=1$: 3 מצבים
    \item $l=2$, $m=-2,-1,0,+1,+2$: 5 מצבים
    \item סה"כ: 9 מצבים
\end{itemize}

\textbf{כלל כללי:}
\[
\text{ניוון רמה } n = \sum_{l=0}^{n-1} (2l+1) = n^2
\]
\end{examplebox}

\begin{theorembox}{ניוון $n^2$ וסימטריה נסתרת}

\textbf{תצפית מפתיעה:} האנרגיה תלויה רק ב-$n$, **לא** ב-$l$!

\textbf{ניוון "צפוי":} $(2l+1)$ — מסימטריה ספירית ($SO(3)$)

\textbf{ניוון "מפתיע":} $n^2$ — גדול בהרבה!

\textbf{מסקנה:} יש **סימטריה נסתרת** שאינה סימטריית סיבוב רגילה.

\textbf{הסבר:} וקטור **\LR{Runge-Lenz}**:
\[
\vec{A} = \frac{1}{2\mu}\left(\vec{p} \times \vec{L} - \vec{L} \times \vec{p}\right) - \frac{Ze^2}{4\pi\epsilon_0}\frac{\vec{r}}{r}
\]

מקיים: $[\hat{H}, \vec{A}] = 0$ (גודל שמור!)

הסימטריה המלאה: $SO(4)$ במרחב המומנטום — חבורה גדולה יותר!

זו תוצאה עמוקה שמסבירה את הניוון הגדול.
\end{theorembox}

\begin{notebox}
\textbf{הערה היסטורית:}

ניוון $n^2$ היה ידוע כבר ל-\LR{Niels Bohr (1913)} ממודל האטום שלו.

\LR{Pauli (1926)} הראה שהסיבה היא וקטור \LR{Runge-Lenz}.

זוהי דוגמה מרתקת לכך ש"מקרה" מתמטי (ניוון גדול) מגלה סימטריה פיזיקלית עמוקה!
\end{notebox}

\subsection{משמעות פיזיקלית: ספקטרוסקופיה}

\subsubsection{מעברים וקווים ספקטרליים}

\begin{importantbox}{מעברים בין רמות}
אלקטרון יכול לקפוץ מרמה $n'$ לרמה $n$ ולפלוט/לספוג פוטון:
\[
h\nu = E_{n'} - E_n = R_\infty Z^2\left(\frac{1}{n^2} - \frac{1}{n'^2}\right)
\]

\textbf{סדרות ספקטרליות:}
\begin{itemize}
    \item **סדרת ליימן** (\LR{Lyman}): $n=1 \leftarrow n'=2,3,4,\ldots$ (\LR{UV})
    \item **סדרת בלמר** (\LR{Balmer}): $n=2 \leftarrow n'=3,4,5,\ldots$ (נראה)
    \item **סדרת פשן** (\LR{Paschen}): $n=3 \leftarrow n'=4,5,6,\ldots$ (\LR{IR})
\end{itemize}
\end{importantbox}

\begin{examplebox}{קו \LR{H-$\alpha$}}
המעבר $3 \to 2$ (\LR{H-$\alpha$}) במימן:
\[
E_\gamma = E_3 - E_2 = 13.6 \, \text{eV} \left(\frac{1}{4} - \frac{1}{9}\right) = 1.89 \, \text{eV}
\]

אורך גל:
\[
\lambda = \frac{hc}{E_\gamma} = 656.3 \, \text{nm} \quad \text{(אדום)}
\]

זהו הקו האדום המפורסם של מימן — נצפה באסטרופיזיקה!
\end{examplebox}

\subsubsection{רוחב הקו ומעברים מותרים}

\begin{notebox}
\textbf{שאלות פתוחות להמשך:}
\begin{enumerate}
    \item האם כל המעברים מותרים? (לא! **כללי ברירה**)
    \item מהו קצב המעבר? (תורת ההפרעות תלויית-זמן)
    \item למה לקווים יש רוחב? (**עקרון אי-וודאות**: $\Delta E \cdot \Delta t \sim \hbar$)
\end{enumerate}

נלמד את כל אלה בהמשך הקורס!
\end{notebox}

\subsection{תיקונים יחסותיים (רמז לעתיד)}

\begin{importantbox}{מבנה דק (\LR{Fine Structure})}
הנוסחה $E_n = -\frac{Z^2}{n^2}R_\infty$ היא **קירוב לא-יחסותי**.

\textbf{תיקונים רלטיביסטיים:}
\begin{itemize}
    \item תיקון קינטי: $\frac{p^4}{8m^3c^2}$
    \item צימוד ספין-מסלול
    \item אפקט דרווין (\LR{Darwin})
\end{itemize}

\textbf{תוצאה:} הניוון ב-$l$ מתבטל חלקית!
\[
E_{nj} = E_n \left[1 + \frac{\alpha^2 Z^2}{n^2}\left(\frac{n}{j+1/2} - \frac{3}{4}\right)\right]
\]

כאשר $j = l \pm 1/2$ (תנע זוויתי כולל).

פרמטר קטן: $\alpha^2 \approx (1/137)^2 \approx 5 \times 10^{-5}$ — אפקט עדין!
\end{importantbox}

\subsection{סיכום כללי — שיעור 07}

\begin{summarybox}

\textbf{1. המשוואה הרדיאלית:}
\begin{itemize}
    \item הפרדנו משתנים: $\psi = R_{nl}(r) Y_{lm}(\theta,\phi)$
    \item תנאי שפה: $r \cdot R(r) \to 0$ כאשר $r \to 0$ (מהרמיטיות $\hat{p}_r$)
    \item החלפה $u = rR$ מפשטת את המשוואה
\end{itemize}

\textbf{2. סקלות אופייניות:}
\begin{itemize}
    \item רדיוס בוהר: $a_0 = \frac{4\pi\epsilon_0\hbar^2}{m_e e^2} \approx 0.529 \, \text{\AA}$
    \item קבוע מבנה דק: $\alpha = \frac{e^2}{4\pi\epsilon_0\hbar c} \approx \frac{1}{137}$
    \item אנרגיית ריידברג: $R_\infty = \frac{\alpha^2 m_e c^2}{2} \approx 13.6 \, \text{eV}$
\end{itemize}

\textbf{3. קוונטיזציה:}
\begin{itemize}
    \item מספר קוונטי ראשי: $n = 1, 2, 3, \ldots$
    \item תנע זוויתי: $l = 0, 1, \ldots, n-1$
    \item הקרנה: $m = -l, \ldots, +l$
\end{itemize}

\textbf{4. ספקטרום האנרגיה:}
\[
E_n = -\frac{Z^2}{n^2} R_\infty = -\frac{13.6 Z^2}{n^2} \, \text{eV}
\]

\textbf{5. ניוון $n^2$:}
\begin{itemize}
    \item רמת $n$ מנוונת $n^2$ פעמים
    \item נובע מסימטריה נסתרת — וקטור Runge-Lenz
    \item חבורת הסימטריה המלאה: $SO(4)$ (לא רק $SO(3)$!)
\end{itemize}

\textbf{6. פיזיקה:}
\begin{itemize}
    \item קווים ספקטרליים: סדרות ליימן, בלמר, פשן
    \item H-$\alpha$: $656.3$ nm (קו אדום של מימן)
    \item תיקונים יחסותיים: מבנה דק $\sim \alpha^2$
\end{itemize}

\end{summarybox}

\subsection*{נקודות למחשבה ובדיקת הבנה}

\begin{enumerate}
    \item הסבירו מדוע $\hat{p}_r = -i\hbar\partial_r$ אינו הרמיטי, ומה צריך להוסיף כדי להפוך אותו להרמיטי.
    
    \item הראו במפורש ש-$\hat{p}_r^2 = -\hbar^2(\partial_r^2 + \frac{2}{r}\partial_r)$ עבור $\hat{p}_r = -i\hbar(\partial_r + \frac{1}{r})$.
    
    \item מדוע תנאי השפה $r \cdot R(r) \to 0$ ב-$r \to 0$ הוא הכרחי? מה יקרה אם לא נדרוש את זה?
    
    \item הסבירו את היתרון בהחלפה $u(r) = r \cdot R(r)$. איך זה משנה את תנאי השפה?
    
    \item חשבו את רדיוס בוהר עבור פוזיטרוניום (אלקטרון+פוזיטרון). למה הוא גדול פי 2 מ-$a_0$?
    
    \item מה המשמעות הפיזיקלית של $\alpha \approx 1/137$? למה קוראים לו "קבוע מבנה דק"?
    
    \item הסבירו מדוע הטור $F(\rho) = \sum a_q \rho^q$ חייב להסתיים. מה היה קורה אם לא?
    
    \item חשבו את אנרגיית היינון של He$^+$ (הליום עם אלקטרון אחד, $Z=2$).
    
    \item מדוע ניוון $(2l+1)$ "מובן" אבל ניוון $n^2$ "מפתיע"? מה זה אומר על הסימטריה?
    
    \item הסבירו מדוע פוטנציאל $V(r) = -\alpha/r^s$ עם $s \neq 1$ לא היה נותן ניוון $n^2$.
    
    \item חשבו את אורך הגל של קו ליימן-$\alpha$ (מעבר $2 \to 1$). באיזו תחום ספקטרלי הוא?
    
    \item מה ההבדל בין ספקטרום הפליטה לספקטרום הבליעה של מימן בטמפרטורת חדר?
    
    \item הסבירו איך ניתן למדוד את $\alpha$ באמצעות ספקטרוסקופיה של אטום מימן.
    
    \item מדוע רמות עם $l$ גדול יותר "יושבות יותר רחוק" מהגרעין? (רמז: פוטנציאל צנטריפוגלי)
    
    \item האם הניוון $n^2$ נשמר כאשר מוסיפים שדה חשמלי חיצוני? (רמז: אפקט שטרק)
\end{enumerate}
